% Compile beside math-review.sty; replace all visible insertion text for publication.
% Rename and regroup functional headings around the actual mathematical thread.
% Shared files do not determine section hierarchy; no environment/count quota applies.
% Edition check: remove a side branch only when the selected explanation stays intact.
% Copyedit: state a full result once locally; check the printed meaning of notation.
\documentclass[12pt,reqno]{amsart}
\usepackage{math-review}
\title[Part IV: Boundaries (concise)]{[Review topic]\\
Part IV: Applications, Boundaries and Problems (concise)}
\author{[Author name]}
\date{[Prepared date; literature cutoff date]}
\begin{document}
\begin{abstract}
\placeholder{Identify the uses and established boundaries of the theory and the core Problems that arise from the remaining mathematical gaps.}
\end{abstract}
\maketitle
\tableofcontents

\section{Introduction}\label{sec:introduction}
\placeholder{State the central question about the theory's reach: what it lets us do, where established boundaries lie and which gaps form the frontier. Explain the mathematical relationship among these themes.}
\placeholder{Explain how the uses and boundaries in Section~\ref{sec:reach} lead to the Problem group in Section~\ref{sec:problems}. Identify the relationship, not merely the names of applications or questions.}

\section{What the theory achieves and where its reach ends}\label{sec:reach}
\placeholder{Introduce only the setting needed for these uses. Choose applications and boundaries for their role in the central question rather than for equal coverage of categories.}

\subsection{A use of the main results}\label{subsec:use}
\placeholder{Verify the hypotheses of an actual result, derive its consequence and explain the significance. A construction or classification can replace a calculation; no particular proof technology is required.}

\subsection{The boundary leading to further questions}\label{subsec:boundary}
\placeholder{Explain a demonstrated limitation or unresolved extension. Do not infer failure of a conclusion from a failed assumption or method. Keep a counterexample, an open range and an evidence gap distinct.}

\section{Core Problems and the remaining scope}\label{sec:problems}
\placeholder{Organize the agreed core frontier by its relationship to the preceding results and boundaries. Give separate environments for genuinely different targets, without a number quota.}
\begin{problem}[A core unresolved target]\label{prob:core}
\placeholder{State the precise objects, assumptions, quantifiers and unresolved goal. Replace this scaffold with an accurately sourced Problem, not a fabricated conjecture.}
\end{problem}
\placeholder{After Problem~\ref{prob:core}, give its formulation source and cutoff status separately from the statement. Explain importance, main verified progress, decisive solved ranges and what remains or could not be verified. Keep this account brief but mathematically specific; label any review-proposed direction.}
\placeholder{For a sharply delimited verified settled scope, replace the open-Problem scaffold with the exact settlement and its sources. A failed search or absent problem list does not justify this exception.}

\templatereferences
\end{document}
